% ============================================================
%  template.tex — Notas Acadêmicas
%  Autor: <Gabriel de Souza Dourado>
% ============================================================

\documentclass[12pt, a4paper]{article}

% ── Codificação e Língua ────────────────────────────────────
\usepackage[utf8]{inputenc}
\usepackage[T1]{fontenc}
\usepackage[brazil]{babel}

% ── Tipografia ──────────────────────────────────────────────
\usepackage{lmodern}
\usepackage{microtype}

% ── Margens ─────────────────────────────────────────────────
\usepackage[
  top=2.5cm, bottom=2.5cm,
  left=2.8cm, right=2.8cm
]{geometry}
\setlength{\skip\footins}{1.5em}

% ── Matemática ──────────────────────────────────────────────
\usepackage{amsmath, amssymb, amsthm}
\usepackage{mathtools}

% ── Algoritmos e Código ─────────────────────────────────────
\usepackage{algorithm}
\usepackage{algpseudocode}
\usepackage{listings}
\usepackage{xcolor}

% ── Pseudocódigo ─────────────────────────────────────
\renewcommand{\algorithmicrequire}{\textbf{Entrada:}}
\renewcommand{\algorithmicensure}{\textbf{Saída:}}

% ── Figuras e Tabelas ───────────────────────────────────────
\usepackage{graphicx}
\usepackage{booktabs}
\usepackage{caption}
\usepackage{float}

% ── Hiperlinks ──────────────────────────────────────────────
\usepackage[
  colorlinks=true,
  linkcolor=blue!70!black,
  citecolor=blue!70!black,
  urlcolor=blue!70!black
]{hyperref}

% ── Cabeçalho e Rodapé ──────────────────────────────────────
\usepackage{fancyhdr}
\pagestyle{fancy}
\fancyhf{}
\fancyhead[L]{\small\textit{\docsubtitle}}
\fancyhead[R]{\small\textit{\docauthor}}
\fancyfoot[C]{\small\thepage}
\renewcommand{\headrulewidth}{0.4pt}

% ── Cores ───────────────────────────────────────────────────
\definecolor{mainblue}{RGB}{30, 80, 160}
\definecolor{codegray}{RGB}{245, 245, 245}
\definecolor{codecomment}{RGB}{100, 100, 100}
\definecolor{codestring}{RGB}{163, 21, 21}
\definecolor{codekeyword}{RGB}{30, 80, 160}

% ── Estilo de Código (C++) ───────────────────────────────────
\lstdefinestyle{cppstyle}{
  language=C++,
  backgroundcolor=\color{codegray},
  basicstyle=\ttfamily\small,
  keywordstyle=\color{codekeyword}\bfseries,
  commentstyle=\color{codecomment}\itshape,
  stringstyle=\color{codestring},
  numberstyle=\tiny\color{codecomment},
  numbers=left,
  numbersep=8pt,
  breaklines=true,
  frame=single,
  framerule=0pt,
  rulecolor=\color{codegray},
  showstringspaces=false,
  tabsize=2,
  captionpos=b,
}
\lstset{style=cppstyle}

% ── Ambientes de Teorema ─────────────────────────────────────
\theoremstyle{definition}
\newtheorem{definition}{Definição}[section]
\newtheorem{example}{Exemplo}[section]

\theoremstyle{plain}
\newtheorem{theorem}{Teorema}[section]
\newtheorem{lemma}[theorem]{Lema}
\newtheorem{corollary}[theorem]{Corolário}
\newtheorem{proposition}[theorem]{Proposição}

\theoremstyle{remark}
\newtheorem*{remark}{Observação}
\newtheorem*{note}{Nota}

% ── Espaçamento ─────────────────────────────────────────────
\usepackage{setspace}
\setstretch{1.3}

% ============================================================
%  METADADOS — altere apenas estas 4 linhas em cada documento
% ============================================================
\newcommand{\docauthor}{Gabriel de Souza Dourado}
\newcommand{\doctitle}{Algoritmos de Machine Learning}
\newcommand{\docsubtitle}{Inteligência Artificial}
\newcommand{\docdate}{\today}
% ============================================================

\begin{document}

% ── Título ───────────────────────────────────────────────────
\begin{center}
  {\large \textcolor{mainblue}{\textsc{\docsubtitle}}}\\[0.4em]
  {\LARGE \textbf{\doctitle}}\\[0.6em]
  {\small \docauthor \quad·\quad \docdate}
  \vspace{0.4em}
  \hrule
\end{center}

\vspace{1em}

% ── Sumário (descomente se quiser) ───────────────────────────
% \tableofcontents
% \newpage

% ------------------------------------------------------------
\section{K-Nearest Neighbors (KNN)}
% ------------------------------------------------------------

O KNN é um algoritmo de aprendizado supervisionado usado tanto para problemas de classificação quanto de regressão. 
É classificado como \textit{Lazy Learning} pois, diferentemente dos algoritmos \textit{eager}, que constroem um modelo 
explícito durante o treinamento, o KNN adia todo o processamento para a fase de predição, limitando-se a memorizar o 
conjunto $\mathcal{D}$.
 
\begin{algorithm}[H]
  \caption{KNN-Treino}
  \begin{algorithmic}[1]
    \Require Conjunto de dados $\mathcal{D} = \{(\mathbf{x}_i, y_i)\}_{i=1}^{N}$,
             onde $\mathbf{x}_i \in \mathbb{R}^{D}$ e $y_i \in \{c_1, \ldots, c_M\}$
    \Ensure  Dados armazenados na memória
    \State \Return $\mathcal{D}$
  \end{algorithmic}
\end{algorithm}
 
\begin{algorithm}[H]
  \caption{KNN-Predição (Classificação)}
  \begin{algorithmic}[1]
    \Require Amostra $\mathbf{z} \in \mathbb{R}^{D}$, inteiro $K \geq 1$,
             conjunto $\mathcal{D}$ armazenado na memória
    \Ensure  Classe predita $\hat{y}$
    \For{$i = 1$ \textbf{até} $N$}
      \State $d_i \gets \|\mathbf{z} - \mathbf{x}_i\|_2$
    \EndFor
    \State $\mathcal{N}_K(\mathbf{z}) \gets$ índices dos $K$ menores valores em $\{d_i\}$
    \State $\hat{y} \gets \underset{c}{\arg\max}
           \displaystyle\sum_{i \in \mathcal{N}_K(\mathbf{z})} \mathbf{1}[y_i = c]$
    \State \Return $\hat{y}$
  \end{algorithmic}
\end{algorithm}

% ------------------------------------------------------------
\section{Decision Trees (Classificação)}
% ------------------------------------------------------------

Consiste em um algoritmo de aprendizado supervisionado usado tanto para problemas de classificação quanto de regressão. 
Formalmente, uma Árvore de Decisão é um \textbf{Grafo Acíclico Direcionado (DAG)} onde cada \textbf{nó interno} 
representa uma condição sobre um atributo, cada \textbf{aresta} representa o resultado dessacondição e cada 
\textbf{folha} armazena a predição final.

\begin{algorithm}[H]
  \caption{ClassTree-MelhorDivisão}
  \begin{algorithmic}[1]
    \Require Subconjunto $\mathcal{D}_m = \{(\mathbf{x}_i, y_i)\}$ no nó $m$,
             critério $Q \in \{\text{Gini},\, \text{Entropia}\}$
    \Ensure  Atributo ótimo $j^*$, ponto de corte ótimo $t^*$ e impureza mínima ponderada
    \State $\text{menor\_imp} \gets \infty$;\; $j^* \gets \text{nulo}$;\; $t^* \gets \text{nulo}$
    \For{$j = 1$ \textbf{até} $D$}
      \State Ordenar $\{x_{ij}\}$ e gerar candidatos $t$ (médias entre valores consecutivos)
      \For{cada candidato $t$}
        \State $R_L \gets \{\mathbf{x} \in \mathcal{D}_m \mid x_{j} \leq t\}$;\;
               $R_R \gets \{\mathbf{x} \in \mathcal{D}_m \mid x_{j} > t\}$
        \State $\text{imp} \gets
               \dfrac{|R_L|}{|\mathcal{D}_m|}\,Q(R_L) +
               \dfrac{|R_R|}{|\mathcal{D}_m|}\,Q(R_R)$
        \If{$\text{imp} < \text{menor\_imp}$}
          \State $\text{menor\_imp} \gets \text{imp}$;\; $j^* \gets j$;\; $t^* \gets t$
        \EndIf
      \EndFor
    \EndFor
    \State \Return $(j^*,\; t^*,\; \text{menor\_imp})$
  \end{algorithmic}
\end{algorithm}

\begin{algorithm}[H]
  \caption{ClassTree-Construção (Recursivo)}
  \begin{algorithmic}[1]
    \Require Subconjunto $\mathcal{D}_m$ e critério $Q$,
    \Ensure  Nó (folha ou interno) da árvore
    \State $\hat{p}_{mk} \gets \dfrac{|\{i : y_i = k\}|}{|\mathcal{D}_m|}$
           \quad para todo $k$
    \If{critério de parada$(\mathcal{D}_m)$ for verdadeiro}
      \State \Return Folha$\!\left(\hat{y}_m \gets \underset{k}{\arg\max}\;\hat{p}_{mk}\right)$
    \EndIf
    \State $(j^*,\; t^*,\; \text{imp\_filhos}) \gets
           \textsc{ClassTree-MelhorDivisão}(\mathcal{D}_m,\; Q)$
    \State $\Delta \gets Q(\mathcal{D}_m) - \text{imp\_filhos}$
    \If{$\Delta < \epsilon$}
      \State \Return Folha$\!\left(\hat{y}_m \gets \underset{k}{\arg\max}\;\hat{p}_{mk}\right)$
    \EndIf
    \State $R_L \gets \{\mathbf{x} \in \mathcal{D}_m \mid x_{j^*} \leq t^*\}$;\;
           $R_R \gets \{\mathbf{x} \in \mathcal{D}_m \mid x_{j^*} > t^*\}$
    \State $\text{filho\_esq} \gets
           \textsc{ClassTree-Construção}(R_L,\; Q)$
    \State $\text{filho\_dir} \gets
           \textsc{ClassTree-Construção}(R_R,\; Q)$
    \State \Return NóInterno$(j^*,\; t^*,\; \text{filho\_esq},\; \text{filho\_dir})$
  \end{algorithmic}
\end{algorithm}

\begin{algorithm}[H]
  \caption{ClassTree-Predição}
  \begin{algorithmic}[1]
    \Require Amostra $\mathbf{z} \in \mathbb{R}^{D}$, árvore $T$
    \Ensure  Classe predita $\hat{y}$
    \State $\text{nó} \gets \text{raiz de } T$
    \While{$\text{nó}$ não é folha}
      \State $(j^*,\; t^*) \gets$ atributo e corte armazenados em nó
      \If{$z_{j^*} \leq t^*$}
        \State $\text{nó} \gets \text{nó.filho\_esq}$
      \Else
        \State $\text{nó} \gets \text{nó.filho\_dir}$
      \EndIf
    \EndWhile
    \State \Return $\hat{y}_m$ armazenado na folha
  \end{algorithmic}
\end{algorithm}

% ============================================================
%  REFERÊNCIAS
% ============================================================

\newpage
\begin{thebibliography}{9}
\bibitem{faceli}
  FACELI, K.; LORENA, A. C.; GAMA, J.; CARVALHO, A. C. P. L. F\@.
  \textit{Inteligência Artificial: Uma Abordagem de Aprendizado de Máquina}.
  Rio de Janeiro: LTC, 2011.

\bibitem{hastie}
  HASTIE, T.; TIBSHIRANI, R.; FRIEDMAN, J\@.
  \textit{The Elements of Statistical Learning: Data Mining, Inference, and Prediction}.
  2.~ed. New York: Springer, 2009.

\bibitem{mitchell}
  MITCHELL, T. M\@.
  \textit{Machine Learning}.
  New York: McGraw-Hill, 1997.

\bibitem{goodfellow}
  GOODFELLOW, I.; BENGIO, Y.; COURVILLE, A\@.
  \textit{Deep Learning}.
  Cambridge: MIT Press, 2016.

\bibitem{bishop}
  BISHOP, C. M\@.
  \textit{Pattern Recognition and Machine Learning}.
  New York: Springer, 2006.
\end{thebibliography}

\end{document}